%% Fonte única da verdade para os metadados da dissertação.
%% Campos ainda desconhecidos permanecem vazios intencionalmente e aparecem
%% como pendências no PDF de rascunho. Não ativar `entrega' em main.tex antes
%% de completar e conferir todos eles.

\titulo{Técnicas de avaliação de sistemas RAG}
\subtitulo{aplicação empírica a discursos parlamentares do Senado Federal}
\titulotraduzido{}

\autor{Fabricio Fernandes Santana}
\autorficha{SANTANA, Fabricio Fernandes}

\orientador{}
\coorientador{}

\curso{Administração Pública}
\programa{Administração Pública}
\area{}
\linhapesquisa{}

\cidade{Brasília}
\local{Brasília-DF}
\data{2027}

\datadefesa{}
% \membrobanca{Nome do orientador}{Orientador}
% \membrobanca[Instituição]{Nome do examinador}{Examinador}

\cutter{}
\numeropaginas{}
\cdd{}
\assuntos{1. Geração aumentada por recuperação. 2. Recuperação da informação. 3. Informação legislativa. 4. Inteligência artificial na administração pública.}

\palavraschave{Geração aumentada por recuperação. Avaliação de sistemas RAG. Informação legislativa. Inteligência artificial. Administração pública.}
\keywords{Retrieval-augmented generation. RAG evaluation. Legislative information. Artificial intelligence. Public administration.}
